\documentclass[10pt,a4paper]{amsart}
\usepackage[margin=19mm]{geometry}
\usepackage{amsmath,amssymb,mathtools}
\usepackage[scr=rsfs,cal=euler]{mathalfa}
\usepackage{xfrac}
\usepackage[unicode,hidelinks,bookmarks=false]{hyperref}
\newcommand{\ie}{i.e.\ }
\newcommand{\cf}{cf.\ }
\newcommand{\resp}{respectively\ }
\newcommand{\Subsection}{\subsection}
\providecommand{\nobreakdash}{\nobreak\hbox{-}}
\setlength{\emergencystretch}{2em}
\multlinegap0pt
\usepackage{amssymb}
\newtheorem{lemma}{Lemma}
\title{Non-convexity of level sets for solutions to $k$-Hessian equations in exterior domains}
\author{Bo Wang, Cong Wang, Zhizhang Wang}
\date{Source: arXiv:2603.28138v1 (CC BY 4.0)}
\begin{document}
\maketitle

\noindent\emph{Source and licence.} Non-convexity of level sets for solutions to $k$-Hessian equations in exterior domains. Version arXiv:2603.28138v1 (CC BY 4.0), distributed by arXiv under the Creative Commons Attribution 4.0 International licence, http://creativecommons.org/licenses/by/4.0/, as recorded in the arXiv licence metadata for this exact version and shown on the abstract page https://arxiv.org/abs/2603.28138v1. Full licence notice: \texttt{licenses/arxiv-2603.28138-CC-BY-4.0.txt}.\par\smallskip

\noindent\textit{Editorial scope.} Counted unit: the article's lemma lem:Phi-k (source0.tex lines 772--774). The article's complete original proof of it is delivered with it (source0.tex lines 775--811, one \texttt{proof} environment of 37 lines). No mathematics is written, repaired or re-derived: the statement and the proof are the article's own lines, in the article's own notation, and no retained line is substituted or re-flowed.  Retained uncounted as the source-local context the proof needs: lines 292--298; lines 763--765.  Nothing was omitted from any retained span, so the package carries no omission record.  No retained line contains a citation, so the wrapper carries no bibliography and no bibliography entry.  No retained line contains a figure environment, an image-inclusion command or drawing code, so no image is delivered and the package declares no asset dependency.  Editorial formatting only, in addition: an AMS \texttt{amsart} preamble that restates the article's own declarations and macros used by the retained lines, namely the environments \texttt{lemma} on separate counters, together with the standard packages those lines need.  The retained environments print in this document as Lemma 1 in order of appearance, whereas the article prints the same items with the numbering of its own full document.  The complete native source of the article as distributed in the arXiv e-print for this exact version is bundled unchanged as \texttt{sources/197411/source0.tex}.
	\begin{equation}\label{eq:subsol}
		\begin{cases}
			S_k(D^2\underline{u}^{\varepsilon,R})\geq 1 & \text{in }\mathbb{R}^n\setminus B_\varepsilon(x_0), \\
			\underline{u}^{\varepsilon,R}=\underline{u}^{\varepsilon} & \text{in }B_2(0)\setminus B_\varepsilon(x_0), \\
			\underline{u}^{\varepsilon,R}=\underline{u} & \text{in }\mathbb{R}^n\setminus E_{2R_0}(0),
		\end{cases}
	\end{equation}

\begin{equation}\label{eq:u-e-v}
	\tilde{u}^\varepsilon\to v\quad\text{in }L^1_{\mathrm{loc}}(\mathbb{R}^n)\cap C^{0,\gamma}_{\mathrm{loc}}(\mathbb{R}^n\setminus\{x_0\}).
\end{equation}

\begin{lemma}\label{lem:Phi-k}
	Let $\tilde{u}^\varepsilon$ and $v$ be defined as above. Then $\tilde{u}^\varepsilon$, $v\in\Phi^k(\mathbb{R}^n)$.
\end{lemma}

\begin{proof}
	It is clear that $\tilde{u}^\varepsilon$ and $v$ are upper semi-continuous and proper in $\mathbb{R}^n$. 

	Let $q$ be a quadratic polynomial such that $\tilde{u}^\varepsilon-q$ has a finite local maximum at some point $\bar{x}\in\mathbb{R}^n  $. If $\bar{x}\in\mathbb{R}^n\setminus\bar{B}_\varepsilon(x_0)$ or $B_\varepsilon(x_0)$, then $S_k(D^2q)=S_k(D^2q(\bar{x}))\geq0$ due to the $k$-convexity and $C^2$ regularity of $\tilde{u}^\varepsilon$ at $\bar{x}$. If $\bar{x}\in\partial B_\varepsilon(x_0)$, we additionally set $\underline{u}^{\varepsilon,R}:=\underline{u}^\varepsilon\leq0$ in $B_\varepsilon(x_0)$. 
	Recalling from \eqref{eq:subsol} that $\underline{u}^{\varepsilon,R}\in C^{\infty}(\mathbb{R}^n\setminus B_\varepsilon(x_0))$ satisfies $\underline{u}^{\varepsilon,R}=\underline u^\varepsilon$ in $B_2(0)\setminus B_\varepsilon(x_0)$, this yields that $\underline{u}^{\varepsilon,R}$ is still smooth and $k$-convex in $\mathbb{R}^n$. Since $u^{\varepsilon,R}$ is increasing with respect to $R$, we have
	$$\tilde{u}^\varepsilon\geq u^{\varepsilon,R}\geq\underline{u}^{\varepsilon,R}\ \text{in }E_R(0)\setminus B_\varepsilon(x_0)\quad\text{and}\quad\tilde{u}^\varepsilon\geq\underline{u}^{\varepsilon,R}\text{ in }B_\varepsilon(x_0).$$ 
	Then for some $\delta>0$ with $B_\delta(\bar{x})\subset E_R(0)  $, 
	$$\tilde{u}^\varepsilon\geq\underline{u}^{\varepsilon,R}\text{ in }B_\delta(\bar{x})\quad\text{and}\quad\tilde{u}^\varepsilon(\bar{x})=\underline{u}^{\varepsilon,R}(\bar{x}).$$
	This yieids $\underline{u}^{\varepsilon,R}-q$ also attains a finite local maximum at $\bar{x}$. Utilizing the $k$-convexity of $C^2$ funciton $\underline{u}^{\varepsilon,R}$, we have $S_k(D^2q)\geq0$. That is, $\tilde{u}^\varepsilon\in\Phi^k(\mathbb{R}^n)$.
	
	We prove $v\in\Phi^k(\mathbb{R}^n)$ by contradiction. If not, there exists a quadratic polynomial $q$, a point $\bar{x}$ and some constant $\delta>0$ such that
	\[v-q\leq0 \text{ in }B_\delta(\bar{x})\quad\text{and}\quad v(\bar{x})-q(\bar{x})=0,\]
	and
	\[S_k(D^2q)<0.\]
	Without loss of generality, we may further assume
	\[
	v - q < c_\delta < 0 \text{ on } \partial B_\delta(\bar{x}). 
	\]
	As otherwise, we can replace the above $q$ with the perturbed polynomial \(q_\beta(x):= q(x) + \beta|x - \bar{x}|^2\) where \(\beta > 0\) is chosen so small that \(S_k(D^2q_\beta) < 0\). By the monotonicity of $\tilde{u}^\varepsilon$ with respect to $\varepsilon$, we obtain 
	\begin{equation}\label{eq:u-q}
		\tilde{u}^\varepsilon - q \leq v - q < c_\delta < 0 \quad \text{on} \, \partial B_\delta(\bar{x}).
	\end{equation} 
	We next consider two cases.

	\textit{Case 1.} If $\bar{x}\neq x_0$, we may assume \(\bar{B}_\delta(\bar{x})\cap\bar{B}_\varepsilon(x_0)=\emptyset\) for sufficiently small \(\varepsilon > 0\). 
	By \eqref{eq:u-e-v}, \(\tilde{u}^\varepsilon(\bar{x}) \to v(\bar{x})\) as $\varepsilon\to 0^+$. That is for any \(\eta > 0\) there exists \(\varepsilon_\eta > 0\) such that when \(0<\varepsilon < \varepsilon_\eta\), \(|\tilde{u}^\varepsilon(\bar{x}) - v(\bar{x})| < \eta\). Taking $\eta=-\frac{c_\delta}{2}$, we get
	\begin{equation*}
		\tilde{u}^{\varepsilon}(\bar{x})>v(\bar{x})-\eta=q(\bar{x})-\eta=q(\bar{x})+\frac{c_\delta}{2}.
	\end{equation*}
	Combining with \eqref{eq:u-q}, \(\tilde{u}^\varepsilon - q\) attains  a finite local maximum in \(B_\delta(\bar{x})\). Since \(S_k(D^2\tilde{u}^\varepsilon) = 1\) in \(B_\delta(\bar{x})\), we have \(S_k(D^2q) \geq 1\), a contradiction.

	\textit{Case 2.} If $\bar{x}=x_0$, we deduce from \( v(x_0) = \limsup_{x \to x_0} v(x) \)   that for any \( \eta > 0 \) small, there exists \(\{x_n\} \subset B_\delta(x_0) \setminus \{x_0\}\) and \( x_n \to x_0 \) such that \(|v(x_n) - v(x_0)| < \frac{\eta}{3}\) and \(|q(x_n)-q(x_0)|<\frac{\eta}{3}\). We fix \( x_n \), then there exists \( \varepsilon_1 = \varepsilon_1(\eta, x_n) > 0 \) such that when \( 0<\varepsilon < \varepsilon_1 \), \( |\tilde{u}^\varepsilon(x_n) - v(x_n)| < \frac{\eta}{3} \). Therefore, for this \( x_n \in B_\delta(x_0)\setminus\{x_0\} \) and $0<\varepsilon<\varepsilon_1$,
	\[
	|\tilde{u}^\varepsilon(x_n) - q(x_n)| \leq |\tilde{u}^\varepsilon(x_n) - v(x_n)| + |v(x_n) - v(x_0)| + |q(x_0) - q(x_n)| < \eta:=-\frac{c_\delta}{2}.
	\]
	Combining with \eqref{eq:u-q}, \( \tilde{u}^\varepsilon - q \) also attains a finite local maximum in \( B_\delta(x_0) \). Since $\tilde{u}^\varepsilon\in\Phi^k(\mathbb{R}^n)$, we have \( S_k(D^2q) \geq 0 \), a contradiction. 
\end{proof}

\end{document}
